%%%%%%%%%%%%%%%%%%%%%%%%%%%%%%%%%%%%%%%%%%%%%%%%%%%%%%%%%%%%%%%%%%%%%%%%% 
% 
% Ejercicios resueltos de Variable Compleja I.
% Doble Grado de Informática y Matemáticas.
% Universidad de Granada.
% Curso 2016/17. 
% 
% 
% Agradecimientos:
% Andrés Herrera (@andreshp) y Mario Román (@M42) por
% las plantillas base.
% 
% Sitio original:
% https://github.com/libreim/apuntesDGIIM/
% 
% Licencia:
% CC BY-NC-SA 4.0 (https://creativecommons.org/licenses/by-nc-sa/4.0/)
% 
%%%%%%%%%%%%%%%%%%%%%%%%%%%%%%%%%%%%%%%%%%%%%%%%%%%%%%%%%%%%%%%%%%%%%%%%% 


% ------------------------------------------------------------------------------
% ACKNOWLEDGMENTS
% ------------------------------------------------------------------------------

%%%%%%%%%%%%%%%%%%%%%%%%%%%%%%%%%%%%%%%%%%%%%%%%%%%%%%%%%%%%%%%%%%%%%%%% 
% Plantilla básica de Latex en Español.
% 
% Autor: Andrés Herrera Poyatos (https://github.com/andreshp) 
% 
% Es una plantilla básica para redactar documentos. Utiliza el paquete fancyhdr
% para darle un estilo moderno pero serio.
% 
% La plantilla se encuentra adaptada al español.
% 
%%%%%%%%%%%%%%%%%%%%%%%%%%%%%%%%%%%%%%%%%%%%%%%%%%%%%%%%%%%%%%%%%%%%%%%%% 

%%%%%%%%%%%%%%%%%%%%%%%%%%%%%%%%%%%%%%%%%%%%%%%%%%%%%%%%%%%%%%%%%%%%%%%%% 
% Plantilla de Trabajo
% Modificación de una plantilla de Latex de Frits Wenneker para adaptarla 
% al castellano y a las necesidades de escribir informática y matemáticas.
% 
% Editada por: Mario Román
% 
% License:
% CC BY-NC-SA 3.0 (http://creativecommons.org/licenses/by-nc-sa/3.0/)
%%%%%%%%%%%%%%%%%%%%%%%%%%%%%%%%%%%%%%%%%%%%%%%%%%%%%%%%%%%%%%%%%%%%%%%%% 

%%%%%%%%%%%%%%%%%%%%%%%%%%%%%%%%%%%%%%%%%%%%%%%%%%%%%%%%%%%%%%%%%%%%%%%%% 
% Short Sectioned Assignment
% LaTeX Template
% Version 1.0 (5/5/12)
% 
% This template has been downloaded from:
% http://www.LaTeXTemplates.com
% 
% Original author:
% Frits Wenneker (http://www.howtotex.com)
% 
% License:
% CC BY-NC-SA 3.0 (http://creativecommons.org/licenses/by-nc-sa/3.0/)
% 
%%%%%%%%%%%%%%%%%%%%%%%%%%%%%%%%%%%%%%%%%%%%%%%%%%%%%%%%%%%%%%%%%%%%%%%%% 


% Tipo de documento y opciones.
%\documentclass[11pt, a4paper, twoside]{article} % Usar para imprimir
\documentclass[11pt, a4paper]{article}


% ---------------------------------------------------------------------------
% PAQUETES
% ---------------------------------------------------------------------------

% Idioma y codificación para Español.
\usepackage[utf8]{inputenc}
\usepackage[spanish, es-tabla, es-lcroman, es-noquoting]{babel}
\selectlanguage{spanish} 
% \usepackage[T1]{fontenc}

% Fuente utilizada.
\usepackage{courier}    % Fuente Courier.
\usepackage{microtype}  % Mejora la letra final de cara al lector.

% Diseño de página.
\usepackage{fancyhdr}   % Utilizado para hacer cabeceras y pies de página.
\usepackage{titlesec} 	% Utilizado para hacer títulos propios.
\usepackage{lastpage}   % Referencia a la última página.
\usepackage{extramarks} % Marcas extras. Utilizado en pie de página y cabecera.
\usepackage[parfill]{parskip}    % Crea una nueva línea entre párrafos.
\usepackage{geometry}            % Geometría de las páginas.

% Símbolos y matemáticas.
\usepackage{amssymb, amsmath, amsthm, amsfonts, amscd}
\usepackage{upgreek}

\usepackage{mdframed}
\usepackage{xcolor}


% Otros.
\usepackage{enumitem}   % Listas mejoradas.
\usepackage[hidelinks]{hyperref}
\usepackage{graphicx}   % Gráficos.
\usepackage{caption}

% Dibujos Tikz
\usepackage{pgf,tikz}
\usepackage{mathrsfs}
\usetikzlibrary{arrows,calc,intersections,through,backgrounds}

% Fuentes personalizadas

\usepackage[scaled=.85]{newpxtext,newpxmath}
\usepackage[scaled=.85]{FiraSans}
\usepackage[T1]{fontenc}

% Código para ajustar las fuentes matemáticas al estilo del texto
% que le rodea

\DeclareMathVersion{sans}
\SetSymbolFont{operators}{sans}{OT1}{cmbr}{m}{n}
\SetSymbolFont{letters}{sans}{OML}{cmbr}{m}{it}
\SetSymbolFont{symbols}{sans}{OMS}{cmbrs}{m}{n}
\SetMathAlphabet{\mathit}{sans}{OT1}{cmbr}{m}{sl}
\SetMathAlphabet{\mathbf}{sans}{OT1}{cmbr}{bx}{n}
\SetMathAlphabet{\mathtt}{sans}{OT1}{cmtl}{m}{n}
\SetSymbolFont{largesymbols}{sans}{OMX}{iwona}{m}{n}

\DeclareMathVersion{boldsans}
\SetSymbolFont{operators}{boldsans}{OT1}{cmbr}{b}{n}
\SetSymbolFont{letters}{boldsans}{OML}{cmbrm}{b}{it}
\SetSymbolFont{symbols}{boldsans}{OMS}{cmbrs}{b}{n}
\SetMathAlphabet{\mathit}{boldsans}{OT1}{cmbr}{b}{sl}
\SetMathAlphabet{\mathbf}{boldsans}{OT1}{cmbr}{bx}{n}
\SetMathAlphabet{\mathtt}{boldsans}{OT1}{cmtl}{b}{n}
\SetSymbolFont{largesymbols}{boldsans}{OMX}{iwona}{bx}{n}

\newif\IfInSansMode
\let\oldsf\sffamily
\renewcommand*{\sffamily}{\oldsf\mathversion{sans}\InSansModetrue}
\let\oldmd\mdseries
\renewcommand*{\mdseries}{\oldmd\IfInSansMode\mathversion{sans}\fi\relax}
\let\oldbf\bfseries
\renewcommand*{\bfseries}{\oldbf\IfInSansMode\mathversion{boldsans}\else%
   \mathversion{bold}\fi\relax}
\let\oldnorm\normalfont
\renewcommand*{\normalfont}{\oldnorm\InSansModefalse\mathversion{normal}}
\let\oldrm\rmfamily
\renewcommand*{\rmfamily}{\oldrm\InSansModefalse\mathversion{normal}}

% Colores

\definecolor{50}{HTML}{E8F5E9}
\definecolor{300}{HTML}{81C784}
\definecolor{500}{HTML}{4CAF50}
\definecolor{700}{HTML}{388E3C}

% ---------------------------------------------------------------------------
% OPCIONES PERSONALIZADAS
% ---------------------------------------------------------------------------

% Redefinir letra griega épsilon.
\let\epsilon\upvarepsilon

% Formato de texto.
\linespread{1.3}            % Espaciado entre líneas.
\setlength\parindent{0pt}   % No indentar el texto por defecto.
\setlist{leftmargin=.5in}   % Indentación para las listas.

% Estilo de página.
\pagestyle{fancy}
\fancyhf{}
\geometry{left=3cm,right=3cm,top=3cm,bottom=3cm}   % Márgenes y cabecera.

% Estilo de las cabeceras

\titleformat{\section}
  {\Large\bfseries\sffamily}{\thesection}{1em}{}
\titleformat{\subsection}
  {\large\sffamily}{\thesubsection}{1em}{}

% Estilo de enlaces
\hypersetup{
  % hidelinks = true,   % Oculta todos los enlaces.
  colorlinks = true,   % Muestra todos los enlaces, sin bordes alrededor.
  linkcolor=black,     % Color de enlaces genéricos
  citecolor={blue!50!black},   % Color de enlaces de referencias
  urlcolor={blue!80!black}     % Color de enlaces de URL
}

% Ruta donde buscar gráficos
\graphicspath{{../Recursos/Plantillas/}, {Recursos/Plantillas/}}

% Redefinir entorno de demostración (reducir espacio superior)
% \makeatletter
% \renewenvironment{proof}[1][\proofname] {\vspace{-15pt}\par\pushQED{\qed}\normalfont\topsep6\p@\@plus6\p@\relax\trivlist\item[\hskip\labelsep\it#1\@addpunct{.}]\ignorespaces}{\popQED\endtrivlist\@endpefalse}
% \makeatother

% Aumentar el tamaño del interlineado
\linespread{1.3}

% Permitir salto de página en ecuaciones
\allowdisplaybreaks

% ---------------------------------------------------------------------------
% COMANDOS PERSONALIZADOS
% ---------------------------------------------------------------------------

% Redefinir letra griega épsilon.
\let\epsilon\upvarepsilon

% Valor absoluto: \abs{}
\providecommand{\abs}[1]{\lvert#1\rvert}    

% Fracción grande: \ddfrac{}{}
\newcommand\ddfrac[2]{\frac{\displaystyle #1}{\displaystyle #2}}

% Texto en negrita en modo matemática: \bm{}
\newcommand{\bm}[1]{\boldsymbol{#1}}

% Línea horizontal.
\newcommand{\horrule}[1]{\rule{\linewidth}{#1}}

% Letras de conjuntos
\newcommand{\R}{\mathbb{R}}
\newcommand{\N}{\mathbb{N}}

% Sucesiones
\newcommand{\xn}{\{x_n\}}
\newcommand{\fn}{\{f_n\}}

% Letra griega "chi" en línea con el texto
\DeclareRobustCommand{\rchi}{{\Large \mathpalette\irchi\relax}}
\newcommand{\irchi}[2]{\raisebox{0.4\depth}{$#1\chi$}} % inner command, used by \rchi 

% Letra 'omega'
\newcommand{\W}{\Omega}
\newcommand{\w}{\omega}


% ---------------------------------------------------------------------------
% CABECERA Y PIE DE PÁGINA
% ---------------------------------------------------------------------------

\renewcommand{\sectionmark}[1]{%
\markboth{#1}{}}

\renewcommand{\subsectionmark}[1]{%
\markright{#1}{}}

%\addtolength{\headheight}{4ex}
\renewcommand{\headrulewidth}{0pt}

\addtolength{\headwidth}{\marginparsep}
%\addtolength{\headwidth}{\marginparwidth}

\fancypagestyle{section}{%
  \fancyhead{}
  %\addtolength{\headheight}{-10ex}
  %\renewcommand{\headheight}{0pt}%
  %\setlength{\footskip}{-48pt}%
  \fancyfoot[LE,RO]{\Large\sffamily\thepage}
  \renewcommand{\headrulewidth}{0pt}%
  \renewcommand{\footrulewidth}{0pt}%
}

\let\originalsection\section
\RenewDocumentCommand{\section}{som}{%
  \IfBooleanTF{#1}
    {\originalsection*{#3}}
    {\IfNoValueTF{#2}
      {\originalsection{#3}}
      {\originalsection[#2]{#3}}%
    }%
  \thispagestyle{section}%
}

\fancyhead[LE,RO]{\rule[-4ex]{0pt}{2ex}\sffamily\textsl{\rightmark}}
\fancyhead[LO,RE]{\sffamily{\leftmark}}
\fancyfoot[LE,RO]{\Large\sffamily\thepage}

% ---------------------------------------------------------------------------
% ENTORNOS PARA MATEMÁTICAS
% ---------------------------------------------------------------------------

% Nuevo estilo para definiciones.
\newtheoremstyle{definition-style} % Nombre del estilo.
{}               % Espacio por encima.
{}               % Espacio por debajo.
{}                   % Fuente del cuerpo.
{}                   % Identación.
{\bf\sffamily}                % Fuente para la cabecera.
{.}                  % Puntuación tras la cabecera.
{.5em}               % Espacio tras la cabecera.
{\thmname{#1}\thmnumber{ #2}\thmnote{ (#3)}}     % Especificación de la cabecera (actual: nombre en negrita).

% Nuevo estilo para notas.
\newtheoremstyle{remark-style} 
{10pt}                
{10pt}                
{}                   
{}                   
{\itshape \sffamily}          
{.}                  
{.5em}               
{}                  

% Nuevo estilo para teoremas y proposiciones.
\newtheoremstyle{theorem-style}
{}                
{}                
{}           
{}                  
{\bfseries \sffamily}             
{.}                
{.5em}           
{\thmname{#1}\thmnumber{ #2}\thmnote{ (#3)}}


% Configuración general de mdframe, los estilos de los teoremas, etc
\mdfsetup{skipabove=12pt,skipbelow=12pt,innertopmargin=12pt,innerbottommargin=4pt}

% Creamos los 'marcos' de los estilos

\mdfdefinestyle{nth-frame}{
	linewidth=2pt, %
	linecolor= 500, % 
	%linecolor=black,
	topline=false, %
	bottomline=false, %
	rightline=false,%
	leftmargin=0pt, %
	innerleftmargin=1em, % 
	rightmargin=0pt, %
	innerrightmargin=0pt, % % 
	%splittopskip=\topskip, %
	%splitbottomskip=\topskip, %
}% 

\mdfdefinestyle{nprop-frame}{
	linewidth=2pt, %
	linecolor= 300, %
	%linecolor= gray, 
	topline=false, %
	bottomline=false, %
	rightline=false,%
	leftmargin=0pt, %
	innerleftmargin=1em, %
	innerrightmargin=0em, 
	rightmargin=0pt, %
	%splittopskip=\topskip, %
	%splitbottomskip=\topskip, %
}%       

\mdfdefinestyle{ndef-frame}{
	linewidth=2pt, %
	linecolor= 300, % 
	%linecolor= gray!50,
	backgroundcolor= 50,
	%backgroundcolor= gray!5,
	topline=false, %
	bottomline=false, %
	rightline=false,%
	leftmargin=0pt, %
	innerleftmargin=1em, %
	innerrightmargin=1em, 
	rightmargin=0pt, % 
	innertopmargin=1.5em,%
	innerbottommargin=1em, % 
	splittopskip=\topskip, %
	%splitbottomskip=\topskip, %
}%               

% Asignamos los marcos a los estilos

\surroundwithmdframed[style=nth-frame]{nth}
\surroundwithmdframed[style=nprop-frame]{nprop}
\surroundwithmdframed[style=nprop-frame]{ncor}
\surroundwithmdframed[style=ndef-frame]{ndef}
\surroundwithmdframed[style=ndef-frame]{ejer}

% Nuevo estilo para ejemplos.
\newtheoremstyle{example-style}
{10pt}                
{10pt}                
{}                  
{}                   
{\bf \sffamily}              
{:}                 
{.5em}               
{}                   

% Teoremas, proposiciones y corolarios.
\theoremstyle{theorem-style}
\newtheorem{nth}{Teorema}[section]
\newtheorem{nprop}{Proposición}[section]
\newtheorem{ncor}{Corolario}[section]
\newtheorem{lema}{Lema}[section]
% Definiciones.
\theoremstyle{definition-style}
\newtheorem{ndef}{Definición}[section]
\newtheorem{ejer}{Ejercicio}[section]

% Notas.
\theoremstyle{remark-style}
\newtheorem*{nota}{Nota}
\newtheorem*{sol}{Solución}


% Ejemplos.
\theoremstyle{example-style}
\newtheorem{ejemplo}{Ejemplo}[section]

% Listas ordenadas con números romanos (i), (ii), etc.
\newenvironment{nlist}
{\begin{enumerate}
    \renewcommand\labelenumi{(\emph{\roman{enumi})}}}
  {\end{enumerate}}

% División por casos con llave a la derecha.
\newenvironment{rcases}
{\left.\begin{aligned}}
    {\end{aligned}\right\rbrace}


% ---------------------------------------------------------------------------
% PÁGINA DE TÍTULO
% ---------------------------------------------------------------------------

% Título del documento.
\newcommand{\subject}{Variable Compleja I: Ejercicios resueltos}

% Autor del documento.
\newcommand{\docauthor}{Doble Grado en Ingeniería Informática y Matemáticas}

% Título
\title{
  \normalfont \normalsize 
  \textsc{Universidad de Granada} \\ [25pt]    % Texto por encima.
  \horrule{0.5pt} \\[0.4cm] % Línea horizontal fina.
  \huge Varible Compleja I\\ \Large{Ejercicios resueltos}\\ % Título.
  \horrule{2pt} \\[0.5cm] % Línea horizontal gruesa.
}

% Autor.
\author{\Large\sffamily{\docauthor}}

% Fecha.
\date{\vspace{-1.5em} \normalsize \sffamily Curso 2016/17}



% ---------------------------------------------------------------------------
% COMIENZO DEL DOCUMENTO
% ---------------------------------------------------------------------------

\begin{document}

\maketitle  % Título.
\vfill
\begin{center}
  \includegraphics{by-nc-sa}  % Licencia.
\end{center}
\newpage
\tableofcontents    % Índice
\newpage




\section{Números complejos}
\begin{ejer}
	Probar que el conjunto de matrices
	$$ M = \left\{ 
	\left( \begin{array}{cc}
	 a & b \\
	-b & a \end{array} \right)
	: \ a, b\in \mathbb{R}\right\} $$
	con las operaciones de suma y producto de matrices, es un cuerpo isomorfo a $\mathbb{C}$.
\end{ejer}


\begin{ejer}
	Calcular la parte real, la parte imaginaria y el módulo de los números complejos
	$$ \frac{i-\sqrt{3}}{1+i}\hspace{1cm}\text{y}\hspace{1cm} \frac{1}{i\sqrt{3}-1} $$
\end{ejer}



\begin{ejer}
	Sea $U=\{ z\in\mathbb{C} : |z|<1 \}$.
	Fijado $a \in U$, se considera la función $f: U\rightarrow\mathbb{C}$ dada por
	$$ f(z) = \frac{z-a}{1-\overline{a}z}\hspace{1cm} \forall z\in U $$
	Probar que $f$ es una biyección de U sobre sí mismo y calcular su inversa.
\end{ejer}

\begin{sol}

$|f(z)| < 1 \ \forall z : |z| <1$
$f^{-1}(z)= \frac{z+a}{1+\overline{a}z}$
Si $|z|=1$, entonces $|f(z)| = |\frac{z-a}{a-\overline{a}z}|$
multiplicando en esta expresión por $\overline{z}$
$ \frac{z-a}{1-\overline{a}z} \overline{z}  =  \frac{1-a\overline{z}}{1-\overline{a}z}$
$f$ tenemos que es holomorfa en el disco, lleva la frontera en la frontera.
\end{sol}


\begin{ejer}
	Dados $z_1, z_2, \ldots , z_n \in C^{\ast}$, encontrar una condición necesaria y suficiente para que se verifique la siguiente igualdad:
	$$ \left|\sum_{k=1}^n z_k\right| = \sum_{k=1}^n |z_k|$$
\end{ejer}


\begin{sol}



Por inducción, todos los números complejos deben tener el mismo argumento, son vectores linealmente dependientes sin que se invierta el signo de ninguno de ellos.

$\exists\lambda_1,...,\lambda_n>0 : \lambda_1 z_1 = \lambda_2 z_2 = ... = \lambda_n z_n$

para probar que es necesaria no hace falta hacer inducción
$ |\sum_{k=1}^n z_n| = |\sum_{k=1}^n \frac{\lambda_1}{\lambda_k} z_1| 
=
|z_1|\sum_{k=1}^n \frac{\lambda_1}{\lambda_k} 
=
|z_1|\sum_{k=1}^n \frac{|z_k|}{|z_1|} 
=
\sum_{k=1}^n |z_k|$

$n=2 \implies |z_1+z_2| = |z_1| + |z_2| \longleftrightarrow z_2 = \lambda z_1$ con $\lambda \in \mathbb{R}$

Sea cierto para $n\in\mathbb{N}$
$|\sum_{k=1}^{n+1} z_k| 
=
|z_{n+1} + \sum_{k=1}^n z_k|
=
|z_{n+1}| + |\sum_{k=1}^n z_k|$

Hemos usado en este último paso que
$|z_{n+1}|+|\sum_{k=1}^n| \geq |\sum_{k=1}^{n+1} z_k|=\sum_{k=1}^{n+1} |z_k| = |z_{n+1}| + \sum_{k=1}^n |z_k|$
donde la otra desigualdad entre los extremos se sabe de la desigualdad triangular.
%Ahora usamos la hipótesis de inducción
$|z_{n+1}| + |\sum_{k=1}^n z_k| 
=
|z_{n+1}| + \sum_{k=1}^n |z_k|$
Por la hipótesis de inducción para $k=2$, tenemos que 
$\exists \lambda >0 : z_{n+1} = \lambda \sum_{k=1}^n z_k$

Nos queda por hacer lo análogo con el resto de números complejos $z_k,k=1...n$
Por la hipótesis de inducción para $k=n, \exists \mu_1,...\mu_n : \mu_1 z_1 = \mu_2 z_2=...= \mu_n z_n$,
como $z_{n+1} = \lambda \sum_{k=1}^n z_k 
=
\lambda \sum_{k=1}^n \frac{\mu_1}{\mu_k} z_1$

$$ |z_1+z_2| = |z_1|+|z_2| 
\Longleftrightarrow
|z_1+z_2|^2 = (|z_1|+|z_2|)^2 
=
|z_1|^2 + 2|z_1||z_2|+|z_2|^2$$
 entonces tenemos
$$z_1 \overline{z_2}=\lambda>0 \text{ y que } z_1 = \frac{\lambda}{|z_2|^2 }z_2 
$$

\end{sol}

\begin{ejer}
	Describir geométricamente los subconjuntos del plano dados por
	$$ A=\{ z\in\mathbb{C} : |z+i|=2|z-i| \} \text{ y } B=\{ z\in\mathbb{C} : |z-i| + |z+i| = 4 \} $$
\end{ejer}

\begin{sol}


Veamos el conjunto A, tenemos en cuenta $z=(a,b)$,
$A$ cumple $ \sqrt{a^2+(b+1)^2} = 2\sqrt{a^2+(b-1)^2}$
, entonces $a^2+b^2+1+2 b = a^2+(b+1)^2 = 4(a^2+(b-1)^2) = 4a^2+4b^2+4-8b$, entonces
$3a^2+3b^2-10b+3=0$, $a^2+b^2-\frac{10}{3}b+1=0$, sumamos y restamos $\frac{25}{9}$, $a^2+(b-\frac{5}{3})^2-\frac{25}{9}+1$, 
$a^2+(b-\frac{5}{3})^2=\frac{16}{9}$
Vemos que A es la circunferencia con $c=(0, \frac{5}{3})$, $r=\frac{4}{3}$

Veamos el B, elevando dos veces al cuadrado tenemos como resultado una elipse
$$ \frac{a^2}{\sqrt{3}^2} + \frac{b^2}{2^2} = 1 $$
\end{sol}



\begin{ejer}
	Probar que $arg(z) = 2\arctan( \frac{Im(z)}{Re(z) + |z|} )$ para todo  $z\in\mathbb{C}^{\ast}\backslash\mathbb{R}^{-}$.
\end{ejer}

\begin{sol}

Hemos usado la fórmula del ángulo doble
$$ \cos(2\alpha) = \cos^2(\alpha)-\sin^2(\alpha) $$
$$ \sin(2\alpha) = 2\sin(\alpha)\cos(\alpha) $$

$\phi = \arctan(\frac{Im z}{Re z + |z|})$
$\phi = 2\alpha = \arctan(\frac{Im z}{Re z + |z|}) $

\textbf{Pista para otra forma de hacerlo}
$$ \phi = 2\arctan ( \frac{Im z}{Re z + |z|} ) 
|z|(\cos(\phi)+i\sin(\phi)) =? z $$
Para ello usamos, si $t=\theta/2$
$$ \cos(\theta) = \frac{1-\tan^2(\theta/2)}{1+\tan^2(\theta/2)}
\sin(\theta) = \frac{2\tan(\theta/2)}{1+\tan^2(\theta/2)} $$
\end{sol}


\begin{ejer}
	Probar que, si $z=x+iy\in\mathbb{C}^{\ast}$, con $x,y\in\mathbb{R}$, se tiene
	
	$$
	arg(z) = \left\{ 
		\begin{array}{lcc}
			\arctan(y/x) &   si  & x>0 \\
			\arctan(y/x) + \pi & si & x<0, y>0 \\
			\arctan(y/x)-\pi & si & x<0,y<0 \\
			\pi/2 & si & x=0,y>0 \\
			-\pi/2 & si & x=0,y<0 
		\end{array}
	\right.
	$$
\end{ejer}


\begin{ejer}
	Probar las \textit{fórmulas de De Moivre}:
	$$ \cos(n\theta) + i\sin(n\theta) = (\cos(\theta)+i\sin(\theta))^n \hspace{1cm}\forall\theta\in\mathbb{R}, \forall n\in\mathbb{N} $$
\end{ejer}


\begin{sol}
Haremos la prueba por inducción, para el caso inicial $n=1$ es trivial, suponemos cierto para un $n\in\mathbb{N}$ genérico.
Probamos que en ese caso es cierto para $n+1$.

$$ 
\cos((n+1)\theta)+i\sin((n+1)\theta) 
=
\cos(n\theta+\theta)+i\sin(n\theta+\theta)
$$
$$ =
\cos(n\theta) \cos(\theta)- \sin(n\theta)\sin(\theta) +i\cos(n\theta)\sin(\theta) + i\sin(n\theta)\cos(\theta) $$
$$ =
\cos(n\theta)(cos(\theta)+i\sin(\theta)) + i\sin(n\theta)(cos(\theta)+i\sin(\theta))
$$
$$=
(\cos(theta)+i\sin(\theta)) (\cos(n\theta)+i\sin(n\theta))
=
(cos(\theta)+i\sin(theta))^{n+1}
$$
La última igualdad la tenemos por la hipótesis de inducción.
\end{sol}



\begin{ejer}
	Calcular las partes real e imaginaria del número complejo 
	$\left( \frac{1+i\sqrt{3}}{2} \right)^8$
\end{ejer}


\begin{ejer}
	Probar que, para todo $x\in\mathbb{R}$, se tiene
	\begin{itemize}
		\item $\sin\left( \frac{x}{2} \right) \sum_{k=0}^{n} \cos(kx) = \cos\left( \frac{nx}{2} \right) \sin\left( \frac{(n+1)x}{2} \right)$
	\end{itemize}
\end{ejer}
\begin{sol}


\textbf{Pista}
Probar las dos simultáneamente y usar la fórmula de Moivre
$ \sum_{k=0}^n ( cos(x)+sin(x) )^k $



$$\sin(x/2) \sum_{k=0}^{n} \cos(kx) + i\sin(x/2) \sum_{k=0}^{n} \sin(kx)$$
Sacamos factor común 
$$ \sin(x/2)( \sum_{k=0}^{n} \cos(kx)+i\sin(kx) ) $$
Usamos la fórmula de Moivre
$$ \sin(x/2) \sum_{k=0}^{n} (\cos(x)+i\sin(x)^k $$
Usamos que $x\not\in 2\pi\mathbb{Z}$
$$ \sin(x/2) \frac{1-(\cos(x)+i\sin(x)^{n+1}}{1-(\cos(x)+i\sin(x)} $$
$$= \sin(x/2) \frac{ 1- (\cos((n+1)x)+i\sin((n+1)x)) }{1 - \cos(x) - i\sin(x)} $$
Usamos que $2\sin^2(x/2) = 1-\cos(x)$ y que $\sin(x) = \sin(x/2)\cos(x/2)$,
$$ 
\sin(x/2) \frac{ 1- (\cos((n+1)x)+i\sin((n+1)x)) }{ 2\sin^2(x/2)-2i\sin(x/2)\cos(x/2) }
=
\frac{ 1- (\cos((n+1)x)+i\sin((n+1)x)) }{ 2\sin(x/2)-2i\cos(x/2) }
 $$
$$
\frac{ 2\sin^2 (\frac{n+1}{x}x)-i2\sin(\frac{n+1}{2}x)\cos(\frac{n+1}{2}x) }{ 2\sin(x/2) -2i\cos(x/2) }
$$
Multiplicamos numerador y denominador por el conjugado del denominador
$$
\sin(\frac{n+1}{2}x)( \sin(\frac{n+1}{2}x)-i\cos(\frac{n+1}{2}x) )( \sin(x/2)+i\cos(x/2) )
$$
$$ 
=  \sin(\frac{n+1}{2}x) ( \cos(\frac{n+1}{2}x)\cos(x/2) + \sin(\frac{n+1}{2}x)\sin(x/2)  + i( \sin(\frac{n+1}{2}x)\cos(x/2) - \cos(\frac{n+1}{2}x)\sin(x/2) ) )
$$
$$ = \sin(\frac{n+1}{2}x) (\cos(\frac{nx}{2}) + i\sin(\frac{nx}{2}))
$$
\end{sol}

\newpage

\section{Topología del plano complejo}
\begin{ejer}
	Estudiar la continuidad de la función argumento principal, $\arg : \mathbb{C}^{\ast} \rightarrow \mathbb{R}$.
\end{ejer}

\begin{sol}


Usamos la fórmula de la relación anterior para probar que es continua en $\mathbb{C}^{\ast} \backslash \mathbb{R}$
$$  
arg z = 2\arctan (\frac{Imz}{Rez + |z|}) \hspace{1cm} \forall z\in\mathbb{C}^{\ast}\backslash\mathbb{R}
$$
Luego nos aproximamos por sucesiones
$$  
\{ z+\frac{i}{n} \} \rightarrow z \hspace{1cm}
\lim_n(arg(z+\frac{i}{n})) = \lim_n( \arctan(\frac{1}{nz}+\pi) ) = \pi
$$

$$
\{ z-\frac{i}{n} \} \rightarrow z
\lim_n(arg(z-\frac{i}{n})) = \lim_n( \arctan(\frac{-1}{nz}-\pi) ) = -\pi
$$
\end{sol}

\begin{ejer}
	Dado $\theta\in\mathbb{R}$, se considera el conjunto $S_{\theta} = \{ z\in\mathbb{C}^{\ast} : \theta\not\in Arg(z) \}$. 
	Probar que existe una función $\varphi\in\mathcal{C}(S_{\theta})$ que verifica $\varphi(z)\in Arg(z)$ para todo $z\in S_{\theta}$
\end{ejer}

\begin{sol}
Definimos
$$
f(z) = z( \cos(\pi-\theta) + i\sin(\pi-\theta) ) \hspace{2cm} f:\mathbb{C}\rightarrow \mathbb{C} \text{ es continua}
$$
$$
\phi(z) = arg(f(z)) - (\pi-\theta) \hspace{2cm} \phi: S_{\theta} \rightarrow \mathbb{C} \text{ $\phi$ es continua}
$$
$$ 
\phi(z) \in Arg (f(z)) + \theta-\pi \subset Arg (f(z)) + Arg( \cos(\theta-\pi)+i\sin(\theta-\pi) ) 
$$
$$= Arg ( f(z)-\cos(\theta-\pi)+i(\theta-\pi) ) = Argz
$$
\end{sol}

\begin{ejer}
	Probar que no existe ninguna función $\phi\in\mathcal{C}^{\ast}$ tal que $\varphi(z)\in Arg(z)$ para todo $z\in\mathcal{C}^{\ast}$, y que el mismo resultado es cierto, sustituyendo $\mathbb{C}^{\ast}$ por $\mathbb{T}=\{ z\in\mathbb{C} : |z|=1 \}$
\end{ejer}

\begin{sol}
$\mathbb{T}$ compacto y conexo $\implies \phi(\mathbb{T})$ intervalo cerrado y acotado

Sea $\alpha\in\mathbb{T} : \phi(\alpha) \not =\min(\phi(\mathbb{T})), \max(\phi(\mathbb{T}))$ 
$\phi ( \mathbb{T}-\{ \alpha \} )$ conexo $\implies [ \min(\phi(\mathbb{T})), \phi(\alpha)[ \cup ]\phi(\alpha), \max( \phi(\mathbb{T}) ) ] $ que es una contradicción.

Si $\exists h:\mathbb{C}^{\ast} \rightarrow \mathbb{R}$ con $h(z)\in Arg(z)$ $\forall z\in\mathbb{C}^{\ast}$ y $h$ continua, entonces
$h_{| \mathbb{T}} : \mathbb{T} \rightarrow \mathbb{T} \rightarrow \mathbb{R}$ es continua y $h_{|\mathbb{T}} (z) \in Arg(z) \ \forall z\in\mathbb{T}$
\end{sol}

\begin{ejer}
	Probar que la función $Arg : \mathbb{C}^{\ast} \rightarrow \mathbb{R}\backslash \mathbb{R}/2\pi\mathbb{Z}$ es continua, considerando en $\mathbb{R}/2\pi\mathbb{Z}$ la topología cociente. Más concretamente, se trata de probar que, si $\{z_n\}$ es una sucesión de números complejos no nulos, tal que $\{z_n\} \rightarrow z\in\mathbb{C}^{\ast}$ y $\theta\in Arg(z)$, se puede elegir $\theta_n\in Arg(z_n)$ para todo $n\in\mathbb{N}$, de forma que $\{ \theta_n \} \rightarrow \theta$.
\end{ejer}




\begin{sol}


$$
\{ z_n \} \rightarrow z \implies \forall \epsilon>0\ \exists m\in\mathbb{N} : n\geq m \ |z_n-z|z\epsilon
$$
Sea $\epsilon = \frac{|z|}{2}$, entonces 
$\exists m\in\mathbb{N} :  n\geq m\ z_n\in D(z, |z|/2)$

A partir de aquí usamos el ejercicio $2$ para concluir.
\end{sol}



\begin{ejer}
	Dado $z\in\mathbb{C}$, probar que la sucesión $\left\{ \left( 1+\frac{z}{n} \right)^n \right\}$ es convergente y calcular su límite.
\end{ejer}


\textbf{Idea}
$$
|1+\frac{z}{n}|^n \rightarrow e^{Re z}
$$
$$
\lim_{n\rightarrow\infty} |1+\frac{z}{n}|^n 
= e^{ \lim_{n\rightarrow\infty} n(|1+\frac{z}{n}|-1) }
$$
$$
arg((1+\frac{z}{n})^n) \rightarrow Im z
$$

\begin{sol}



$z_n\in\mathbb{C}, \phi_n\in Arg(z_n)$

$\{ |z_n| \} \rightarrow |z|$ y $\{y_n\} \rightarrow y\in Arg(z)$

Vamos a usar que is tiene sun sucesión de numeros complejos y la sucesion de los modulos converge y hay una sucesion de los argumentos de forma q convergen, la sucesion converge al modulo por el cos +isen

$\phi_n \in Arg(z_n)$, $z_n = (1+\frac{z}{n})^n = (1+\frac{a+ib}{n})^n
= [(1+\frac{a}{n}) + (\frac{ib}{n})] ^n$
$|z_n| = [| (1+\frac{a}{n}) + \frac{ib}{n} |] ^n
= [ (1+\frac{a}{n})^2 + (\frac{b}{n})^n ]^{1/2}
= lim n [ (1+\frac{a^2}{n^2} + (\frac{a}{n}+\frac{b^2}{n^2}) )]^{n/2}
= e^{  \lim n/2 (\frac{a^2}{n^2} + 2a/n + b^2/n^2 )  } = e^a = e^{Re z}$ 

Vamos a utilizar la fórmula de Moivre.

$ z_m = (1+z/n)^n$
$Arg(a+ \frac{a+ib}{n})^n = n Arg(a+\frac{a+ib}{n}) = n*\arctan (\frac{b/n}{1+a/n}) = n*\arctan(\frac{b}{n+a})$
Llamamos $y=b/(n+a)$ y usamos $\lim_{y\rightarrow0} \frac{\arctan(y)}{y} = 1$,
$n*\arctan(\frac{b}{n+a}) = \frac{\arctan(\frac{b}{n+a})}{ \frac{n+a}{n}\frac{b}{n+a}\frac{1}{b} } = b$
\end{sol}
\newpage

\section{Funciones holomorfas}
\subsection{Ejercicio 1}
\subsubsection{a)}
Cauchy Riemman

\textbf{Solución}


\textbf{c)}

$f(x+iy) = \frac{x^3+iy^3}{x^2+y^2}$, $\forall (x,y)\in\mathbf{R}^2 \backslash \{ (0,0) \}$, $f(0)=0$

$u(x,y) = Re f = \frac{x^3}{x^2+y^2}, v(x,y) = Imf = \frac{y^3}{x^2+y^2}$

En $\mathbf{C}^{\ast}$

$ \frac{\partial u}{\partial x} = \frac{3x^2}{x^2+y^2} + \frac{x^4(-1)2}{(x^2+y^2)^2}$
$ \frac{\partial v}{\partial y} = \frac{3y^2}{x^2+y^2} = \frac{y^4(-1)2}{(x^2+y^2)^2}$

entonces
$\frac{\partial u}{\partial y} = \frac{x^3 (-1)2y}{(x^2+y^2)^2}$
$\frac{\partial v}{\partial x} = \frac{y^3(-1)2x}{(x^2+y^2)^2}$

De las dos ecuaciones primeras deducimos

$3x^2(x²+y^2) - 2x^4 = 3y^2(x^2+y^2)-2y^4$  (1)

de últimas dos deducimos

$2yx^3=-2xy^3 \implies xy(x^2+y^2) = 0 \implies xy=0 ó x^2+y^2=0$

por tanto

$x^4-y^4 = 0 \implies |x|=|y|$



\subsection{Ejercicio 2}
2 cauchy riemman
$Imf = 4x^3*y - 4xy^3$

\subsection{Ejercicio 3}
3 cauchy rieman
$a=-c$


\subsection{Ejercicio 4}
\textbf{Enunciado}


\textbf{Solución}
\textbf{4 cauchy rieman, derivar respecto x y respecto y}
$Ref (x,y) = u(x,y)$ $Imf(x,y) = v(x,y)$

Los casos extremos son $a=0$ o $b=0$, donde o la imagen o la parte real que quedan serían constantes.
$au(x,y)+bv(x,y)=c$
$-bv(x,y) = c-au(x,y)$
$v(x,y)= c'-a'u(x,y)$
Derivamos

$\frac{\partial v(x,y)}{\partial x} = \frac{-a' \partial u(x,y)}{\partial x} = \frac{-a' \partial v(x,y)}{\partial y}$
$= \frac{\partial v(x,y)}{\partial y} = \frac{-a' \partial u(x,y)}{\partial y} = -a' \frac{\partial v(x,y)}{\partial x}$
implica
$\frac{\partial u (x,y)}{\partial x} = -a'\frac{\partial v(x,y)}{\partial y} = -a'a'\frac{\partial v(x,y)}{\partial x}$
por tanto
$\frac{\partial v(x,y)}{\partial x} = 0 = \frac{\partial v(x,y)}{\partial y}$ en $\Omega$


\subsection{Ejercicio 5}
\textbf{Enunciado}

Sea $\Omega$ un dominio y $f\in\mathcal{H}(\Omega)$. Probar que si $\overline{f}\in\mathcal{H}(\Omega)$, entonces $f$ es constante.

\textbf{Solución}

$\Omega$ dominio, $f\in\mathcal{H}(C)$
$f+\overline{f} = 2Re f$
$g$ es holomorfa por ser suma de holomorfas, $g\in\mathcal{H}(C)$
$g\in\mathcal{H}(\Omega)$ y su parte imaginaria es constante, por tanto $g$ es constante, entonces

$Re f$ es constante, $f\in\mathbb{H}(\Omega), \Omega$ dominio implican que $f$ es constante.


\subsection{Ejercicio 6}
\textbf{Enunciado}





\textbf{Solución}

$P(z) = a_n z^n + a_{n-1}z^{n-1}+...+a_1z + a_0$
$P(\overline{z}) = a_n \overline{z}^n + a_{n-1}\overline{z}^{n-1}+...+a_1\overline{z} + a_0 $

$\overline{ P(\overline{z}) } = \overline{a_n}z^n + \overline{a_{n-1}}z^{n-1}+...+\overline{a_1}z + \overline{a_0}$

$\overline{z}^n  = \overline{z} ... \overline{z} = \overline{(z^n)}$


resolvemos

$a\in\Omega^{\ast} \longleftrightarrow \overline{a}\in\Omega$
$\lim_{z\rightarrow a} \frac{f^{\ast}(z)-f^{\ast}(a)}{z-a} = \lim_{z\rightarrow a} \frac{\overline{f(\overline{z})} - \overline{f(\overline{a})}}{z-a} =
 \lim_{z\rightarrow a} \frac{\overline{\overline{f(\overline{z})} - \overline{f(\overline{a})}}}{z-a} =
 \lim_{z\rightarrow a} \overline{\left( \frac{f(\overline{z})-f(\overline{a})}{z-a} \right)} =
 \overline{f'(\overline{a})}$


\subsection{Ejercicio 7}
\textbf{Enunciado}



\textbf{Solución}

\textbf{Idea} 7 jugar con los grados del polinomio

Suponemos que  existe una función racional que es la exponencial:

$R(z) = \frac{\sum_{i=0}^k \lambda_i z^i}{\sum_{i=0}^m \mu_i z^i} = \frac{A(z)}{B(z)}$

$R(z) = R'(z) = \frac{A'(z)B(z)-A(z)B'(z)}{B(z)^2} = \frac{A(z)}{B(z)}$ 
(Tachamos los $B(z)$)

$A(z) = A'(z) - \frac{B'(z)A(z)}{B(z)}$,
$B(z)A(z) = A'(z)B(z) - B'(z)A(z)$, esto se da $\forall z\in D(a,r)$ con $a\in\mathbb{C}, r>0$

Tienes dos polinomios que son iguales en un entorno no vació. Para deducir lo último tienes que pensar que son polinomios iguales en un montón de puntos.


%$gr(A(z)) \leq gr(A'(z))$
$gr(B(z)) + gr(A(z)) \leq \max \{ gr(A'(z)) + gr(B(z)), gr(A(z))+gr(B'(z)) \}$

Por tanto llegamos a una contradicción, nuestra hipótesis no era cierto.
\newpage

\section{Funciones analíticas}
\subsection{Ejercicio 1}
\subsubsection{b)}


\textbf{Solución}

$\sum_{n\geq 0} z^{2n} = \sum_{n\geq 0} \alpha_nz^n$
con $\alpha_nz{2n+1} =0, \alpha_{2n} = 1$
hay que ver el límite superior de la sucesión
$\limsup |\alpha_n| = 1 \implies R=\frac{1}{1} = 1$

\subsubsection{c)}

\textbf{Solución}

$\sum_{n\geq 0} 2^nz^{n!}$,
si $|z|\geq 1$, entonces la serie diverge, 
si $|z|<1$
$\sqrt[n]{|2^nz^{n!}|} = \sqrt[n]{2^n|z|^{n!}} = 2*|z|^{(n-1)!} \rightarrow 0$
por tanto
$\sum_{n\geq 0} |2^nz^{n!}| < \infty \implies \sum_{n\geq 0} 2^nz^{n!} < \infty$, entonces $R=1$.

\subsubsection{e)}

\textbf{Solución}

$\sum_{n\geq 0} (n+a^n)z^n$ con $a\in\mathbb{R}^+$.
Lo separamos en dos series, calculamos los radios de convergencia por separado y cogemos el menor.

$\sum_{n\geq 0} nz^n$ y $\sum_{n\geq 0} a^nz^n$

% ---------------------- hay que arreglarlo
$a^n \leq n+a^n \leq (n+1)a^n$
en los extremos el radio de convergencia es $1/a$, por tanto el radio de convergencia de $n+a^n$ es $1/a$

$\sum_{n\geq 0} a^n |z|^n \leq \sum_{n\geq 0} (n+a^n)|z|^n \leq \sum_{n\geq 0} (n+1)a^n|z|^n$

Por tanto el radio de convergencia es $1/a$...................................
% ----------------------

//Podemos usar el criterio de la raíz para sucesiones 
$\sum_{n\geq 0} (n+a^n)z^n$ con $a\in\mathbb{R}^+$
Calculamos $\lim_{n\rightarrow\infty} \sqrt[n]{n+a^n}$
$ = \lim_{n\rightarrow\infty} \frac{(n+1)+a^{n+1}}{n+a^n} = a (\text{si a>1}), 1 (\text{ si a<=1})$
por tanto
$R = 1/a (a>1), 1 (a<=1)$

\subsubsection{c)}
$\sum_{n\geq 0} a^{n^2}z^n,a\in\mathbb{C}$
raíz n-esima 


\subsection{Ejercicio 2}
\subsubsection{a)}
No se reduce el radio de convergencia, se puede intentar el radio de convergencia de las derivadas


\subsection{Ejercicio 3}
Convergen en todo el plano las que a partir de cierto término son 0, o sea, las que son una suma finita

Sea $\sum_{n\geq 0} \alpha_n (z-a)^n$ que converge uniformemente,
$S_n=\sum_{k=0}^{n-1} \alpha_k (z-a)^k$ c.u. $\Longleftrightarrow S_n$ es uniformemente de Cauchy $\Longleftrightarrow \forall\epsilon>0,\exists m\in\mathbb{N}: p\geq q\geq n \implies |S_p(z)-S_q(z)| < \epsilon \forall z\in\mathbb{C}$
Vemos que los polinomios divergen en infinito.
Sea $p\in\mathbb{P}[\mathbb{C}]$, $\lim_{|z|\rightarrow\infty} p(z) = \infty \Longleftrightarrow gr(p)>0 $

Vemos que $|S_p(z)-S_q(z)|$ es un polinomio que está acotado, tiene que ser constante para no diverger en infinito, en particular es la constante $0$, ya que $S_p(a)=S_q(a)$, por tanto $\sum_{q=m}^{p-10} \alpha_n (z-\alpha)^n = 0 \implies  \alpha_n=0 \forall n\geq m$

\subsection{Ejercicio 4}
Idea:

llamar w a (w-1)/(z-1), y mirar cuando la función tiene imagen que cae en el disco de centro 0 y radio 1, afinar para saber cuando un subconjunto de C po esta función cae dentro de un subconjunto compacto donde el sumatorio converja

\newpage

\section{Funciones elementales}
\subsection{Ejercicio 1}


\textbf{Solución}

$f(0) = f(0+0) = f(0)^2$
$f(0)= 0$ ó $f(0) = 1$

Si $f(0)=0$ la función es constante.
EN el caso $f(0)\not =0$, $f(0)=1$, 
Si $f$ es derivable en $\alpha\in\mathbb{C}$

$\exists \lim_{n\rightarrow 0} \frac{f(\alpha+h)-f(\alpha)}{n}$
vemos que cumple la fórmula de adición 
$\frac{f(\alpha+h)-f(\alpha)}{n} = f(\alpha) \frac{f(n)-f(0)}{n} \Longleftrightarrow f$ es derivable en $0$
Y eso se puede aplicar $\forall z\in\Omega$

Ahora encontramos todas las funciones enteras que cumplan la condición del enunciado.

Sea $z\in\Omega$, $f(z)=0$, $f(z) = e^{wz} : w\in\mathbf{C}$

Sea $f$ tal que $f\in\mathbb{H}(\mathbb{C})$

$ f'(z) = f(z) \lim_{h\rightarrow 0} \frac{f(h)-f(0)}{h}$
$(f(z) e^{-wz})' = f'(z) e^{-w}z + (-w) +ze^{-wz} = f(z)we^{-wz} - f(z)we^{-wz} = 0$
las funciones son la misma salvo una constante
En el punto $0$ las dos funciones valen lo mismo, por lo que la función es la constante $1$, ya que $f(0)=1$


El ejemplo es $f(z) = e^{Re(z)}$


\subsection{Ejercicio 2}

$e^z = e^{Re(z) (\cos(Imz) + i\sin(Imz))}$
$B_V = \{ z\in\mathbb{C} : a\leq Rez \leq b \}$, $a,b\in\mathbb{R}$ con $a<b$
$B_H = \{ z\in\mathbb{C} : a\leq Imz \leq b \}$
si $a\leq Rez \leq b$, $e^a \leq e^{Rez} \leq e^b$
lo que pasa es que se puede mover en toda la circunferencia unidad

cuando la parte imaginaria se puede mover donde quiera le das infinitas vueltas a la circunferencia unidad. Tenemos que

$e^B_V$ e la corona circular de centro $0$ y radios $e^a$ y $e^b$

donde

$a\leq Imz \leq b$

Tenemos que $\exp(B_H)$ es el sector del plano encerrado entre los ángulos $a$ y $b$




\subsection{Ejercicio 4}
\textbf{Enunciado}

Probar que si $\{z_n\}$ y $\{w_n\}$ son sucesiones de números complejos, con $z_n \not = 0$ para todo
$n\in \mathbb{N}$ y $\{ z_n \} \rightarrow 1$, entonces

$$ \{w_n(z_n-1)\} \rightarrow \lambda\in\mathbb{C} \implies \{ z_m^{w_n} \}\rightarrow e^{\lambda} $$

\textbf{Solución}
Como la función exponencial es continua
$\{ w_n (z_n-1) \} \rightarrow \lambda \implies \{ e^{w_n(z_n-1)} \} \rightarrow e^{\lambda}$

$\lim \{ \log(z_n) \} = 0 \implies \lim \{ z_n-\log(z_n) \} = 1$
$\implies \lim \{ w_n(z_n-\log(z_n))-w_n \} = 0$

$\lim \{ z_n^{w_n} = e^{\log(z_n)-w_n} \}= \lim \{ e^{w_n (z_n-1)} \}$
$\Longleftrightarrow \lim\{ \frac{e^{w_n(z_n-1)}}{e^{w_n (\log(z_n))}} = e^{w_n (z_n-\log(z_n))-w_n} \} = 1$

Vemos que
$z_n^{w_n} = e^{w_n \frac{\log(z_n)}{z_n-1} (z_n-1)}$
$= e^{w_n (z_n-1) \frac{\log(z_n)}{z_n-1}}$
Sabemos que $ \frac{\log(z_n)}{z_n-1} \rightarrow 1$ ya que
$\lim_{z\rightarrow 1} \frac{\log(z)-\log(1)}{z-1} = \log'(1) = 1/1$





\subsection{Ejercicio 5}
\textbf{Enunciado}

Estudiar la convergencia puntual, absoluta y uniforme de la serie de funciones
$\sum_{n\geq 0} e^{-nz^2}$

\textbf{Solución}

La serie converge puntualmente si, y sólo si, $| \frac{1}{e^{z^2}} | <1 \Longleftrightarrow 1<|e^{z^2}|  \Longleftrightarrow 0<Rez^2 = (Rez)^2-(Imz)^2 $
$\Longleftrightarrow |Rez| > |Imz|$

donde en la última implicación hemos usado
$e^{z^2} = e^{Rez^2} e^{Imz^2}$

$Rez^2 = (Rez)^2 - (Imz)^2$,
$Rez^2>0 \Longleftrightarrow |Rez| > |Imz|$


Vemos ahora la convergencia uniforme

$A = \{ z\in\mathbb{C} : (Rez)^2>(Imz)^2 \}$
Si $B\subset A$ y satisface que $\inf_{z\in B} [ (Rez)^2-(Imz)^2 ] > 0$, entonces hay convergencia uniforme en $B$




\subsection{Ejercicio 6} % ----
\textbf{Enunciado}

Probar que $a, b, c \in\mathbb{T}$ son vértices de un triángulo equilátero si, y sólo si, $a+b+c = 0$ .


\textbf{Solución}

$\{ a,b,c \} = \{ e^(\lambda - 2/3\pi), e^{\lambda i}, e^{(\lambda + 2/3\pi)i} \}$

$\Rightarrow$

$e^{2/3\pi i} (a+b+c) = e^{2/3\pi i} a + e^{2/3\pi i}b + e^{2/3\pi i}c = a+b+c \implies a+b+c=0$

$\Leftarrow$

$a'=\frac{a}{a} = 1, b'=\frac{b}{a}, c'=\frac{c}{a}$
$b'=e^{\theta i}$, $c' = e^{\gamma i} = -b-1$


$a+b+c = 0 \implies a'+b'+c' = 0 \implies 1+b'+c' = 0 \implies c'=-b'-1$ con $\gamma,\theta\in ]-\pi,\pi[$

De lo que deducimos que 
$(-\cos(\theta)-1)-i\sin(\theta) = (\cos(\gamma)) + i(\sin(\gamma))$

$-\in(\theta) = \sin(\gamma) \implies \theta = -\gamma $
$-\cos(\theta) -1 =  \cos(\gamma)$
por tanto
$\theta = \pm\frac{2\pi}{3} = -\gamma$



\subsection{Ejercicio 7}




\textbf{Solución}

$a\in\Omega$ y vemos que es derivable por la definición
$\lim_{z\rightarrow a} \frac{\phi(z)-\phi(a)}{z-a} \frac{\phi(z)-\phi(a)}{\phi(z)+\phi(a)} = \lim_{z\rightarrow a} \frac{z-a}{z-a} \frac{1}{\phi(z)+\phi(a)} = \frac{1}{2\phi(a)} = \phi '(a)$
donde hemos usado que $\phi(a)\not = 0$




\subsection{Ejercicio 8}
\textbf{Enunciado}

Probar que, para todo $z \in D(0, 1)$ se tiene:

a) $\sum_{n=1}{\infty} \frac{(-1)^{n+1}}{n} z^n = \log(1+z)$


\textbf{Solución}
a)
$\log (1+z) \in\mathbb{H}(D(0,1))$ y $(\log (1+z))' = \frac{1}{1+z}$ $\forall z\in D(0,1)$

$\frac{1}{1+z} = \frac{1}{1-(-z)} = \sum_{n=0}^{\infty} (-1)^n z^n$
por otra parte la serie de potencias 

$\sum_{n\geq 1} \frac{(-1)^{n+1}}{n} z^n$ tiene radio de convergencia $1$
y su suma $\sum_{n=1}^{\infty} \frac{(-1)^{n+1}}{n} z^n = f(z)$ es holomorfa en $D(0,1)$ y su derivada se calcula término a término

$f'(z) = \sum_{n=1}^{\infty} \frac{(-1)^{n+1}}{n} n z^{n-1} = \sum_{n=0}^{\infty} (-1)^{n}z^{n}$

Entonces $f'(z) = g'(z)$ $\forall z\in D(0,1)$, por tanto $f$ y $g$ difieren en una constante.

Como $g(0) = \log(1) = 0 = f(0)$
con lo que tenemos que $f$ y $g$ son iguales en $D(0,1)$



\subsection{Ejercicio 9}


\textbf{Solución}

$f(z) = \log(\frac{1+z}{1-z})$
la función es holomorfa en $\mathbb{C}\backslash \{1\}$
sabemos que $\log \in \mathbb{H}(\mathbb{C}^{\ast}\backslash\mathbb{R}^-)$
y vemos cuando la función logaritmo cae dentro de dicho conjunto

de forma intuitiva
$\frac{1+z}{1-z} \in\mathbb{R}^- \Longleftrightarrow \exists r>0 : z\not=1, \frac{1+z}{1-z} = -r \Longleftrightarrow 1+z = rz-r \Longleftrightarrow z(r-1)=1+r \Longleftrightarrow z = \frac{1+r}{r-1}$

Viendo que
$g(z) = \frac{1+z}{1-z} \in\mathbb{H}(\Omega)$ y $g(\Omega) \subseteq \mathbb{C}^{\ast}\backslash \mathbb{R}^-$
podemos asegurar que $f\in\mathbb{H}(\mathbb{C})$ por composición.

$f'(z) = \frac{ \frac{1-z+(1+z)}{(1-z)^2} }{ \frac{1+z}{1-z} } = \frac{2}{1-z^2}$

Siendo $\xi\in\mathbb{H}(\Omega)$, entonces
$\xi '(z) = \frac{1}{1+z} + \frac{1}{1-z} = \frac{2}{1-z^2}$



\textbf{Pista}
hacer $\frac{2}{1-z^2}$ en serie de potencias





\subsection{Ejercicio 10}
\textbf{Pista}
$z^{\phi} = e^{ \phi \log z}$ con $\Omega,\Omega_{\phi} \subset \mathbb{C}^{\ast}\backslash \mathbb{R^-}$

$f^{-1}(z) = \phi^{\frac{1}{\phi\log z}} = z^{1/{phi}}$


\subsection{Ejericicio 12}
\textbf{Pista}

$\sin(nz) = \frac{e^{inz}-e^{-inz}}{2i}$, entonces
$\frac{1}{2i} \sum_{n\geq 0} \frac{e^{inz}-e^{-inz}}{2^n}$
Tenemos que ver cuando, para un $z\in\mathbb{C}$ fijo, estudiar la convergencia de las series
$\sum_{n\geq 0} \frac{e^{inz}}{2^n}$ y $\sum_{n\geq 0} \frac{e^{-inz}}{2^n}$

$e^{inz} = e^{in(Rez+iImz)} = e^{-nImz+inRez} = e^{-nImz}e^{inRez}$ con $|e^{inRez}|=1$

$|\frac{e^{inz}}{2^n}| = \frac{e^{-nImz}}{2^n} = (\frac{e^{-Imz}}{2})^n$
entonces

$\sum_{n\geq 0} |\frac{e^{inz}}{2^n}|$ converge $\Longleftrightarrow e^{-Imz} < 2 \Longleftrightarrow -Imz < \ln 2 \Longleftrightarrow -\ln 2 < Imz$

y tenemos convergencia uniforme en $B\subset A= \{ z\in\mathbb{C} : |Imz|<\ln 2 \}$ tal que $\sup_{z\in B} |Imz| < \ln 2$
\newpage

\section{Integral curvilínea}
\subsection{Ejercicio 1}
\textbf{Enunciado}

\textbf{Solución}

$\gamma :[0,r] \rightarrow \mathbb{C}$,
$\gamma (s) = \alpha + s, s\in[0,r]$

$\int_{[\alpha,\alpha + r]} f(z)dz = \int_{0}^{r} f(\alpha +s)\gamma ' (s)ds$, donde $\gamma ' (s) = 1 \forall s\in [0,r]$

$\int_{0}^{r} f(\alpha +s)\gamma ' (s)ds = \int_{0}^{r} f(\alpha + s)ds$

$\xi : [0,r] \rightarrow \mathbb{C}$, $\xi(s) = \alpha +is \forall s\in[0,r]$

$\int_{[\alpha,\alpha + ir]} f(z)dz = \int_0^r f(\alpha + is) \gamma ' (s)ds = i\int_0^r f(\alpha +is)ds$



\subsection{Ejercicio 2}
\textbf{Enunciado}

\textbf{Solución}
 
Tenemos que probar que
$\lim_{r\rightarrow\infty} \int_{\gamma_r} \frac{zdz}{z^3+1} = 0$

$| \int_{\gamma_r} \frac{z}{z^3+1}dz | \leq \int_{\gamma_r} \frac{|z|}{|z^3+1|} dz \leq l(\gamma_r)M(r)$
donde $M(r)>0$ satisface que $\frac{|z|}{|z^3+1|} \leq M(r) \forall z\in\gamma_r{\ast}$

Dado $z\in\mathbb{C}(0,r)^{\ast}$, $|z| = r $

$\frac{|z|}{|z^3+1|} \leq \frac{r}{|z^3|-1} = \frac{r}{r^3-1}$

Por tanto
$l(\gamma_r)M(r) = 2\pi r \frac{r}{r^3-1}$ y $\lim_{r\rightarrow\infty} \frac{r}{r^3-1} = 0$

\

Para la otra integral tenemos que probar que

$\lim_{r\rightarrow\infty} \int_{\gamma_r} \frac{z^2 e^z}{z+1} dz = 0$ 

$| \int_{\gamma_r} \frac{z^2 e^z}{z+1}dz | < \int_{\gamma_r} \frac{|z|^2 |e^z|}{z+1}dz$

Si $z\in\gamma_r^{\ast}$
$$ |z|^2 = |-r+is|^2 \leq (r+1)^2 \hspace{1cm} s\in[0,1] $$
$$ |z+1| \geq r-1 $$

Como $z\in\gamma_r^{\ast}$, $|e^z| = e^{Rez} = e^{-r}$




\subsection{Ejercicio 3}
\textbf{Enunciado}

\textbf{Solución}

$\int_{\mathbb{C}(0,r)} \frac{\log(i+z)}{z} dz$

donde $0<r<1$, $\mathbb{C}(0,r)^{\ast} \subset D(0,1)$ definimos $f(z) = \log(1+z)$, $f\in\mathcal{H}(D(0,1))$

Usando la fórmula de Cauchy para la circunferencia:

$\int_{\mathbb{C}(0,r)} \frac{\log(i+z)}{z} dz = f(0)2\pi i = 0$

\

Sin usar la fórmula de Cauchy para la circunferencia vemos:

$\log (1+z) = \sum_{n=1}^{\infty} (-1)^{n+1} \frac{z^n}{n}$,
la serie $\sum_{n\geq 1} (-1)^{n+1} \frac{z^n}{n}$ converge uniformemente sobre compactos de $D(0,1)$.

La sucesión $\{ \frac{1}{z} \}$ está acotada en el compacto $\mathbb{C}(0,1)^{\ast}$ 

Usando \textbf{observación} 
$\frac{1}{z} \sum_{n\geq 1} (-1)^{n+1} \frac{z^n}{n}$ converge uniformemente en $\mathbb{C}(0,r)^{\ast}$ a $\frac{\log (1+z)}{z}$

\textbf{Observación}

$\{ f_n \} $ converge uniformemente en $B$ a $f$  
y $g:B\rightarrow\mathbb{C}$ está acotada en $B$, entonces
$\{ gf_n \}$ converge uniformemente en $B$ a $gf$

($\exists M>0 : |g(z)| \leq M \hspace{0.5cm} \forall z\in B  $)


---------------------

$0 = \int_{\mathbb{C}(0,r)} \frac{\log(1+z)}{z} dz = \int_{-\pi}^{\pi}  \frac{\log(1+re^{it})}{re^{it}} dt$
$= i\int_{-pi}^{\pi} \log(1+re^{it}) dt$
$i\int_{-\pi}^{\pi} \ln|1+re^{it}|dt -\int_{-\pi}^{\pi} arg(1+re^{it})dt$
$\implies \int_{-\pi}^{\pi} \ln |1+re^{it}|dt = 0$

Sabemos que $\log(1+re^{it}) = \ln|1+re^{it}| + i*arg(1+re^{it})$
y que
$\gamma(t) = re^{it}, \gamma:[-\pi,\pi] \rightarrow \mathbb{C}$, $\gamma '(t) = ire^{it}$


también sabemos que
$|1+re^{it}| = ((1+r\cos(t))^2 + r^2\sin^2(t))^{1/2} = (1+2r\cos(t) + r^2\cos^2(t) + r^2\sin^2(t))^{1/2} = (1+2r\cos(t)+r^2)^{1/2}$


por tanto
$0 = \int_{-\pi}^{\pi} \ln|1+re^{it}| dt = 1/2 \int_{-\pi}^{\pi} \ln(1++2r\cos(t) +r^2)dt$

Como sabemos que el coseno es par
$ 1/2 \int_{-\pi}^{\pi} \ln(1++2r\cos(t) +r^2)dt = \int_{0}^{\pi} \ln (1+2r\cos(t)+r^2) dt$


\subsection{Ejercicio 4}
\textbf{Enunciado}

\textbf{Solución}

Sabemos que $f\in D(0,1)$, teniendo que $|f(z)-1|<1, \forall z\in D(0,1)$

Sabemos que $f$ no se anula en $D(0,1)$ y deducimos que
$f(D(0,1)) \subset D(1,1)$.
Por tanto $\log(f) \in \mathcal{H}(D(0,1))$

$\log(f(z))' = \frac{f'(z)}{f(z)} \implies \frac{f'}{f}$ admite primitiva holomorfa en $D(0,1)$
y $\mathbb{C}(0,r) \subset D(0,1)$ es un camino cerrado, por tanto
$\int_{\mathbb{C}(0,r)} \frac{f'(z)}{f(z)}dz = 0$



\subsection{Ejercicio 5}
\textbf{Enunciado}

\textbf{Solución}

$\frac{1}{1+z^2} = \frac{1}{(z-i)(z+i)} = \frac{A}{z-i} + \frac{B}{z+i} = \frac{1}{2i} \left[ \frac{1}{z-i} - \frac{1}{z+i} \right]$

Por tanto

$\int_{\mathbb{C}(i,1)} \frac{1}{1+z^2} dz = \frac{1}{2i} \left[ \int_{\mathbb{C}(i,1)}\frac{1}{z-i}dz - \int_{\mathbb{C}(i,1)} \frac{1}{z+i}dz \right] \not = 0$

La primera es $2\pi i$

La segunda integral es $0$ por el teorema de Cauchy para dominios estrellados, teniendo que $\mathbb{C}(i,1) \subset D(i,3/2)$.



\subsection{Ejercicio 6}
\textbf{Idea}

$\Omega = \mathbb{C}\backslash \{ iy : |y|>1 \}$
$\arctan \in \mathcal{H}(\Omega)$ y $\arctan$ es una primitiva de $\frac{1}{1+z^2}$

$\sigma ^{\ast} \subset \Omega$ y $\sigma$ es un camino cerrado.

Con esa información sabremos que 
$\int_{\sigma} \frac{1}{1+z^2} dz = 0$


\newpage


\section{Teorema local de Cauchy}
\subsection{Ejercicio 1}


\textbf{Solución}

Fijo $z\in\mathbf{C}$ con $|z-a|>r, \exists r'>0 : |z-a|>r'>r$
Consideramos $\Omega = D(a,r'), f(w) = \frac{1}{w-z}$ , $f\in\mathcal{H}(D(a,r'))$
$D(a,r')$ es convexo, en particular estrellado, por tanto, usando $TLC$ y $TCDE$ $\int_{C(a,r)} \frac{dw}{w-z} = 0$

\subsection{Ejercicio 2}


\textbf{Solución}

Fijo $z\in D(a,R)$
Fijo $0<r<R : |z-a|<r<R$, $\overline{D}(a,r) \subset D(a,R)$

Por la fórmula de Cauchy para la circunferencia $f(z) = \frac{1}{2\pi i} \int_{C(a,r)} \frac{f(w)}{w-z} dw$

$ \left|\frac{1}{2\pi i} \int_{C(0,R)} \frac{f(w)}{w-z} dz - \frac{1}{2\pi i} \int_{C(0,r)} \frac{f(w)}{w-z}dw \right| = \frac{1}{2\pi} |\int_{-\pi}^{\pi} \frac{f(a+Re^{it}) iRe^{it} }{ a+Re^{it}-z } dt - \int_{-\pi}^{\pi} \frac{f(a+re^{it}) ire^{it}}{a+re^{it}-z} dt|$

$\leq \frac{1}{2\pi} \int_{-\pi}^{\pi} \left| \frac{f(a+Re^{it})R}{a+Re^{it}-z} -\frac{f(a+re^{it})r}{a+re^{it}-z} \right| dt$

$g:[r,R]x]-\pi, \pi[ \rightarrow \mathbb{C}$,
$g(p,t) = \frac{f(a+pe^{it})p}{a+pe^{it}-z}$
por lo que el término de la integral queda $|g(R,t)-g(r,t)|$

$g$ es una función continua en un compacto, por lo tanto es uniformemente continua.

Dado $\epsilon >0 \exists \delta : si \| (p,t)-(p',t') \| _{\infty}<\delta \implies |g(p,t)-g(p',t')| < \epsilon$,
y lo utilizamos para $(p,t)=(R,t)$ y $(p',t')=r,t$, si $|R-r|<\delta$ entonces $|g(R,t)-g(r,t)| < \epsilon$

Por lo tanto

$\frac{1}{2\pi} \int_{-\pi}^{\pi} \left| \frac{f(a+Re^{it})R}{a+Re^{it}-z} -\frac{f(a+re^{it})r}{a+re^{it}-z} \right| dt \leq \frac{1}{2\pi}\epsilon 2\pi$ en el caso de que $|R-r|<\delta$


\subsection{Ejercicio 3}

\textbf{Solución}

$\frac{1}{(z-b)(z-c)} = \frac{1}{b-c} [\frac{1}{z-b}-\frac{1}{z-c}]$

Vemos primero el caso $b\not = c$

Tenemos tres subcasos
\begin{itemize}
	\item b interior, c exterior $\int_{C(a,r)} f(z)dz = \frac{1}{b-c} [ \int_{C(a,r)} \frac{dz}{z-b} - \int_{C(a,r)} \frac{dz}{z-c} ] = \frac{1}{b-c} (2\pi i (b interior) - 0 (ej1))$
	
	\item b y c interiores $\int_{C(a,r)} f(z)dz =  \frac{1}{b-c}[2\pi i - 2\pi i] = 0$
	
	\item b y c exteriores, por el ejercicio 1 ambas partes son $0$: $\int_{C(a,r)} f(z)dz = 0-0 = 0$
\end{itemize}

Ahora vemos el caso $b=c$
$f(z) = \frac{1}{(z-b)^2}$ tiene primitiva holomorfa en $\mathbb{C}^{\ast}, \frac{-1}{z-b} \implies \int_{C(a,r)} f(z)dz = 0 $


\subsection{Ejercicio 4}

\textbf{Solución}
\subsubsection{a}

$\int_{C(0,r)} \frac{z+1}{z(z^2+4)} dz$, donde $r\not = 2$, $z\not = \pm 2i$

$\frac{z+1}{z(z^2+4)} = \frac{A}{z} + \frac{B}{z-2i} + \frac{C}{z+2i}$, donde $A = 1/4, B=1/8-i/4, C = -1/8-i/4$

Por tanto

$\int_{C(0,r)} \frac{z+1}{z(z^2+4)} dz = \frac{1}{4} \int_{C(0,r)} \frac{dz}{z} + (\frac{1}{8}-\frac{i}{4})\int_{C(0,r)} \frac{1}{z-2i}dz + (-\frac{1}{8}+\frac{i}{4}) \int_{C(0,r)} \frac{1}{z+2i}$

para $r\not = 2$ las dos últimas integrales se anulan, y por tanto $\frac{1}{4}\int_{C(0,r)} \frac{dz}{z} = \frac{2\pi i}{4} = \frac{\pi}{2}$



\subsection{Ejercicio 5}

\textbf{Solución}

$R>\max\{ |a|, |b| \}$

$\int_{C(0,R)} \frac{f(z)}{(z-a)(z-b)}dz = \frac{1}{b-a} \left[ \int_{C(0,R)} \frac{f(z)}{z-b}dz - \int_{C(0,R)} \frac{f(z)}{z-a}dz \right]$

Por el ejercicio 3 tenemos

$\frac{f(z)}{(z-a)(z-b)} = \frac{f(z)}{b-a} \left[ \frac{1}{z-b} - \frac{1}{z-a} \right]$

Por la fórmula de Cauchy para la circunferencia

$= \frac{1}{b-a} \left[ \int_{C(0,R)} \frac{f(z)}{z-b}dz - \int_{C(0,R)} \frac{f(z)}{z-a}dz \right] = 2\pi i f(b) - 2\pi if(a)$

\textbf{Consecuencia}

Si $\exists M>0 : |f(z)|\leq M \forall z \in \mathbb{C}$

$\frac{2\pi |f(b)-f(a)|}{|b-a|} = \frac{|2\pi i(f(b)-f(a))|}{|b-a|} = \left| \int_{C(0,R)} \frac{f(z)}{(z-a)(z-b)} dz \right|$
$\leq \int_{C(0,R)} \frac{|f(z)|}{|z-a||z-b|} dz \leq \frac{M}{(R-|a|)(R-|b|)} 2\pi R$ que tiende a $0$ cuando $R\rightarrow \infty$ ya que $|z-a|\geq R-|a|$ y $|z-b|\geq R - |b|$
\newpage

\section{Equivalencia entre analiticidad y holomorfía}
\subsection{Ejercicio 1}

\textbf{Solución}

$\phi (w,z) = \frac{\varphi}{w-z}$, 

$\phi : \gamma^{\ast}x\mathbb{C}\backslash \gamma^{\ast} \rightarrow \mathbb{C}$

$\phi$ es continua

$\phi_w(z) = \phi(w,z)$

$z\rightarrow \phi_w(z)$

$\phi_w \in \mathcal{H}(\mathbb{C}\backslash \{w\})$, $\phi_w \in \mathcal{H}(\mathbb{C} \backslash \gamma^{\ast}  )$

Por tanto  por el teorema de holomorfía de una integral dependiente de $1$ parámetro $f\in\mathcal{H}(\mathbb{C}\backslash \gamma^{\ast})$

Podemos proceder también de otra forma:


Fijo $a\in\mathbb{C}\backslash \gamma^{\ast}$
$\lim_{z\rightarrow a} \frac{f(z)-f(a)}{z-a}$
$= \lim_{z\rightarrow a}  \frac{ \int_{\gamma} \frac{\varphi (w)}{w-z}dw  -  \int_{\gamma} \frac{\varphi (w)}{w-a}dw }{z-a}$

$= \lim_{z\rightarrow a} \int_{\gamma} \varphi \frac{(w-a)-(w-z)}{(z-a)(w-z)(w-a)} dw$
$= \lim_{z\rightarrow a} \int_{\gamma} \varphi (w)\frac{dw}{(w-z)(w-a)}$

Sabemos que podemos intercambiar límite e integral

Nos podemos restringir a un compacto para ver $\psi : \phi_{| \gamma{\ast}  \times  \overline{D}(a,r)} \rightarrow \mathbb{C}$  es continua con $r$ suficientemente pequeño y $\psi (w,z) = \frac{\varphi (w)}{(w-z)(w-a)}$

tenemos que $\gamma^{\ast} \times \overline{D}(a,r)$ es compacto, por tanto $\phi$ es uniformemente continua.

La fórmula para las derivadas la podemos obtener por inducción:

$k=1$ lo hemos hecho, supuesto cierto para $k$, vemos para $k+1$



%$\frac{\phi (w)}{(w-z)^{k+1}} - \frac{\phi (w)}{(w-a)^{k+1}} = %\phi (w) = \left[ \frac{(w-a)^{k+1} - %(w-z)^{k+1}}{(w-z)^{k+1}(w-a)^{k+1}} \right]$






\subsection{Ejercicio 2}

\textbf{Solución}

$(1+z)^{\alpha} = e^{ \alpha \log(1+z) } \in\mathcal{H}(D(0,1))$ por tanto es analítica en ese disco.

Por el Teorema de desarrollo de Taylor
$\sum_{n\geq 0} \frac{f^{n)}(0)}{n!}z^n$


\subsection{Ejercicio 3}

\textbf{Solución}

\subsubsection{a} 
$f'(z) = \frac{2z-3}{z^2-3z+2}$

$\frac{1}{z^2-3z+2} = \frac{A}{2-z} + \frac{B}{1-z}$

Por tanto el desarrollo en serie de potencias de la original se puede conseguir a partir de desarrollo de las dos partes por separado.

$f'(z) = (2z-3) \sum_{n\geq 0} (\frac{A}{2} (\frac{z}{2})^n + Bz^n) = (2z-3)\sum_{n\geq 0} (\frac{A}{2^{n+1}} + B)z^n $

$= \sum_{n\geq 0} (\frac{A}{2^{n}} + 2B)z^{n+1} - \sum_{n\geq 0} 3 (\frac{A}{2^{n+1}}+B)z^n$

que nos da como resultado 

$$ f'(z) = -\left(\frac{3A}{2}+B\right) + \sum_{n\geq 1} \left(\frac{A}{2^{n+1}} - B \right)z^n $$


\subsection{Ejercicio 6}

\textbf{Solución}

$f(z) = \frac{1}{1-z-z^2} = \frac{1}{(z-a_1)(z-a_2)} $
$z^2+z-1 = 0 \Longleftrightarrow z = \frac{-1\pm \sqrt{5}}{2}$, donde $a_1=\frac{-1+ \sqrt{5}}{2},a_2=\frac{-1- \sqrt{5}}{2}$
$f(z) = \frac{1}{(a_1-z)(a_2-z)} = \frac{1}{a_1-a_2}\left[ \frac{1}{a_1-z}-\frac{1}{a_2-z} \right]$
Expresamos $\alpha_n$ en términos de $a_1$ y $a_2$
$a_1+a_2 = -1$
$a_1*a_2=-1$
$(z-a_1)(z-a_2) = z^2-(a_1+a_2) + a_1a_2$

$\alpha_{n+2}$ a partir de aquí se puede expresar en términos de $\alpha_n$ y $\alpha_{n+1}$


\subsection{Ejercicio 1}

\textbf{Solución}
\subsubsection{a)}
$f(z) = \sum_{n\geq 0} \frac{f^{n)}(0)}{n!} z^n$
$ f^{n)}(0) = n$, vemos $\left\{ \frac{n}{n!} \right\}$,  $\left\{ \sqrt[n]{\frac{n}{n!}} \right\} \rightarrow 0 \implies R = \infty$,


\subsubsection{b)} 
$\frac{f^{n)}(0)}{n!} = \frac{(n+1)!}{n!} = n+1$, $\sqrt[n]{n+1} \implies R=1$
Por tanto no hay función entera que cumpla la condición

\subsubsection{c)}
$\frac{f^{n)}(0)}{n!} = \frac{2^n n!}{n!} = 2^n$, $\sqrt[n]{z^n} = 2 \implies R=1/2$

\subsubsection{d)}
$\frac{f^{n)}(0)}{n!} = \frac{n^n}{n!}$, $\sqrt[n]{ \frac{n^n}{n!} }$ por el criterio del cociente
$\implies R=\frac{1}{e} < \frac{1}{2}$



\subsection{Ejercicio 12}

\textbf{Solución}

$\int_{\mathcal{C}(0,2)} \frac{dz}{z^2(z-1)^2}$
$f(z) = \frac{1}{z^2(z-1)^2}$

$\phi(z) = \frac{1}{z}$

$\gamma{\ast} = \mathcal{C}(0,1/2)^{\ast}$
$\phi(\gamma^{\ast}) = \mathcal{C}(0,2)^{\ast}$
Si $\gamma(t) = \frac{1}{2}e^{-it}$
$\phi(\gamma(t)) = 2e^{it}\implies \phi o \gamma = \mathcal{C}(0,2)$ 

Por tanto

$\int_{\mathcal{C}(0,2)} \frac{dz}{z^2(z-1)^2} = \int_{\mathcal{C}(0,1/2)} \frac{1}{ \frac{1}{w^2} (\frac{1}{w}-1)^2} \frac{1}{w^2} dw =$
$\int_{\mathcal{C}(0,1/2)} \frac{w^2dw}{(1-w)^2}$
\newpage

\end{document}
